\section{A coherent additive character for the global Fourier step}
\label{sec:additive-normalization}

The local Fourier computations of MA-CQ use a specified character
and its self-dual measure. Their shell supports and absolute Gauss
normalizations hold for either additive sign. A global Poisson
identity has the extra requirement that the character annihilate
the diagonal arithmetic lattice. That requirement must be checked
before the character PW construction.

\paragraph{MA-NR-001 (the printed signs and their explicit repair).}
Read the finite character in Masullo Example~A.3.1.6, p.~557,
with the natural inclusion
$\mathbb Q_p/\mathbb Z_p=\mathbb Z[1/p]/\mathbb Z\to\mathbb R/\mathbb Z$
printed there. It is
$\psi_p^+(x)=\exp(2\pi i\{x\}_p)$, while
Example~A.3.1.5 on the same page gives
$\psi_\infty^+(x)=\exp(2\pi ix)$.
Their product is the character defined in Example~A.3.1.8,
p.~558. With these literal signs, the self-annihilator assertion
for $\mathcal O_{F,S}$ in Example~A.3.1.11, p.~559,
following (A.3.1.10), is false already for
$F=\mathbb Q$, $S=\{\infty,3\}$.
In fact, for $\Gamma=\mathbb Z[1/3]$ and
$x=1$, $y=1/3$, both diagonally embedded in $\Gamma$,
\[
 \psi_\infty^+(xy)\psi_3^+(xy)
                 =\exp(4\pi i/3)\ne1.
 \tag{MA-NR.1}
\]
Thus even the asserted inclusion
$\Gamma\subseteq\Gamma^\perp$ fails under the printed pairing.
This is a sign convention defect, not a defect of the printed
self-dual Haar measures.

A coherent choice, used for the global constructions here, is
\[
 \vartheta_\infty(x)=\exp(2\pi ix),\qquad
 \vartheta_p(x)=\exp(-2\pi i\{x\}_p),\qquad
 \widehat f(y)=\int f(x)\overline{\vartheta(xy)}\,d\nu(x).
 \tag{MA-NR.2}
\]
For number fields compose these characters with the local traces;
at a complex place the archimedean character is
$\exp(2\pi i\operatorname{Tr}_{\mathbb C/\mathbb R}z)$.
The finite Gauss constants in MA-CQ and MA-CGP are defined by
their displayed integrals with this choice. Changing a finite
additive sign reflects the corresponding Fourier coordinate.
It preserves the finite norm-support bound and the value at zero,
so the corrected source PW spaces are unchanged by this repair.

\begin{proof}
The displayed counterexample follows because the ordinary
$3$-adic fractional part of $1/3$ is $1/3$.
For the repaired global rational character, every rational $r$
satisfies
\[
                    r-\sum_p\{r\}_p\in\mathbb Z.
 \tag{MA-NR.3}
\]
Only primes in its denominator contribute. To prove the assertion,
subtract the indicated finite sum. At each finite prime $p$ the
$p$-th summand removes the negative-power part of $r$, while the
other summands are $p$-integral. The rational difference is
therefore integral at every prime, and hence an integer.
Equation (MA-NR.3) proves that the product in (MA-NR.2)
is trivial on the diagonal $\mathbb Q$.

The exact local repair can also be checked without invoking a
global duality theorem. For any prime $p$, use the pairing
\[
 (x_\infty,x_p),(y_\infty,y_p)\longmapsto
    \exp\bigl(2\pi i(x_\infty y_\infty-\{x_py_p\}_p)\bigr)
 \tag{MA-NR.4}
\]
on $\mathbb R\times\mathbb Q_p$.
The diagonal $\Gamma=\mathbb Z[1/p]$ annihilates itself,
by (MA-NR.3). Conversely suppose $(x_\infty,x_p)$
annihilates $\Gamma$. Testing against $1$ gives
$x_\infty-\{x_p\}_p\in\mathbb Z$, so
$x_\infty=r\in\Gamma$.
Subtract the diagonal point $(r,r)$, which already annihilates
$\Gamma$. Testing the remaining point $(0,x_p-r)$ against
$p^{-k}$ for every $k\ge0$ gives
$x_p-r\in p^k\mathbb Z_p$ for every $k$.
Their intersection is zero, so $x_p=r$.
This proves $\Gamma^\perp=\Gamma$ for the repaired pairing.
Equivalently, under the two positive signs the annihilator is
$\{(r,-r):r\in\mathbb Z[1/p]\}$, which explains the exact
coordinate reflection missing from the printed claim.

If a Fourier kernel $\overline{\psi(xy)}$ is replaced by
$\overline{\psi(-xy)}$, its transform is $\widehat f(-y)$.
Reflection preserves Haar measure, absolute norm, and the origin.
Thus the self-dual measure and all norm-support conditions remain
the same. The local annihilator of $\mathcal O_v$ is still the
inverse different, since it is closed under negation. The local
Gauss phase must instead be recomputed using the chosen kernel;
MA-CQ defines it by an integral, so no unverified phase convention
is imported. These observations prove all asserted repair claims.
\end{proof}

\paragraph{Scope and trace-framework placement.}
This finding uses neither $A$, $M$, nor zero location.
It does not refute the self-duality of the adelic measure in
Example~A.3.1.8: the product with the printed signs still gives
an isomorphism with the Pontryagin dual. What fails is the
identification of the diagonal lattice with its annihilator,
used in the arithmetic Poisson step. The corrected construction
uses a coherent character throughout; the choice is recorded
rather than absorbed into an unnamed root-number scalar.

The multiplicative finite-monodromy characters and their
inertia-invariant Euler polynomials, placed in SGA~1,
Expos\'e~V, \S7 and Deligne, Definition~6.1 and
\S\S6.2--6.3 in MA-CH, do not change under an additive
Fourier-sign repair. This is a normalization of the analytic
Fourier pairing and of the associated Gauss constants, not a
new motive or weight assertion. In the cote~1 trace comparison
(scans~079--080 and~111--112 via BR-TR-01/02), the paired
operators must use the same corrected Fourier identification
before trace cyclicity is applied.
